% BANKS-SOURCE: pdf=198 print=188 kind=section id=9.12
\subsection{Renormalization-group equations for masses and the hierarchy problem}

% BANKS-SOURCE: pdf=198 print=188 kind=source-page
% BANKS-SOURCE: pdf=198 print=188 kind=prose
So far we have talked mostly about the use of RG equations for Euclidean Green
functions. In Lorentzian signature, various complications arise, which we will
sketch in this section. First consider the RG equation for a connected
two-point function
% BANKS-SOURCE: pdf=198 print=188 kind=equation
\[
  \mathcal{D}W_2=-2\gamma W_2,
\]
where $\mathcal{D}$ is the infinitesimal RG operator
$(\mu\partial_\mu+\beta\partial_\alpha+m\gamma_m\partial_m)$ in spinor QED
and $\gamma$ is the anomalous dimension of the field in question (this
argument is applicable to any operator that undergoes multiplicative
renormalization). Since the RG equation is independent of momentum, it applies
near a pole, where
% BANKS-SOURCE: pdf=198 print=188 kind=equation
\[
  W_2\sim\frac{Z}{p^2+M^2-\ii0}.
\]
Since $\mathcal{D}$ is a first-order differential operator, it produces both a
single-pole and a double-pole term when acting on $W_2$. The two cannot cancel
out, and hence must satisfy the equation separately. Thus
% BANKS-SOURCE: pdf=198 print=188 kind=equation
\[
  \mathcal{D}Z=-2\gamma Z
\]

% BANKS-SOURCE: pdf=199 print=189 kind=source-page
and
% BANKS-SOURCE: pdf=199 print=189 kind=equation
\[
  \mathcal{D}M^2=0.
\]
The latter equation is particularly interesting when the theory has no
relevant operators, as is true for a chirally symmetric abelian or non-abelian
gauge theory, like massless QED. In massless QED, we get the equation
% BANKS-SOURCE: pdf=199 print=189 kind=equation
\[
  \left(\mu\frac{\partial}{\partial\mu}+\beta\partial_\alpha\right)M^2=0.
\]
On combining this with dimensional analysis we have the solution
% BANKS-SOURCE: pdf=199 print=189 kind=equation
\[
  M=\mu\exp\!\left(-\int^\alpha\frac{\dd x}{\beta(x)}\right)
  \longrightarrow \mu c\,\ee^{3\pi/(2\alpha)}.
\]
This equation makes no sense unless $c=0$. If $c\ne0$ the theory would not be
continuous at $\alpha=0$, contradicting the assumption of perturbation theory
on which we based the calculation.

This is a very general result: a marginally irrelevant perturbation cannot
produce a finite mass scale. On the other hand, if the perturbation were
marginally relevant, the sign of the exponent would change, and we would learn
that exponentially small masses could be generated. This result contains the
seed of one of the most important ideas in quantum field theory. It gives us an
idea of how the vast hierarchy of mass scales in the real world might be
generated from a system with a single intrinsic scale as large as the Planck
mass. All that is necessary is that the theory have a low-energy effective
approximation that is a marginally relevant perturbation of a fixed-point
quantum field theory. In that case, if the underlying high-energy theory can
give us an explanation of why the initial values of parameters are moderately
close to the critical surface, we can explain the natural occurrence of scales
exponentially smaller than the Planck scale.

The nice behavior of marginally relevant parameters is to be contrasted with
what happens if we have truly relevant parameters, like scalar masses near the
Gaussian fixed point. In that case, the RG tells us that we need a fine tuning
of initial conditions with accuracy $(m/M_P)^2$ in order to explain a mass
scale of order $m$. Most physicists consider this unnatural. By contrast, the
tuning of the bare coupling of QCD at a Planck scale cut-off, which is required
in order to explain the QCD scale in terms of the Planck scale, is only about
one part in $25$.\footnote[17]{I am here assuming that the particle content of the
minimal supersymmetric standard model is the only thing that renormalizes the
QCD coupling between laboratory scales and the Planck scale. Extra matter makes
the tuning of the bare coupling even less severe, as long as it preserves
asymptotic freedom.} Fermion masses, although they seem like relevant
perturbations, actually behave like marginal perturbations because of chiral
symmetry. This means that, if the underlying theory has enough symmetry to
forbid the appearance of a fermion mass in the high-energy Lagrangian, then we
could explain a small value of this mass in the real world in terms of
spontaneous breaking of this symmetry by the IR effects of a marginally
relevant coupling.

% BANKS-SOURCE: pdf=200 print=190 kind=source-page
The concept of technical naturalness is based on this example. We have a
parameter that could be classed as relevant, but might naturally be small as a
consequence of a fundamental symmetry broken spontaneously by a marginally
relevant coupling. It is then said that a theory containing a small value of
this parameter is technically natural, because there will not be power-law
renormalizations of its value. It becomes truly natural when we supply the
underlying explanation for the symmetry, and the dynamical symmetry-breaking
mechanism.

The standard model contains one parameter, the quadratic term in the Higgs
potential, that is not technically natural. It receives quadratically
divergent renormalizations, and it is not easy to understand why it is not of
order the Planck scale. This is called the gauge hierarchy problem, since this
parameter determines the scale of the massive weak gauge bosons. Proposals to
solve this problem divide roughly into two classes. The first goes under the
name of technicolor, and replaces the Higgs field by a new QCD-like sector with
a dynamical scale of order a few TeV. This gives an elegant account of gauge
symmetry breaking, but suffers from numerous phenomenological problems when we
try to couple it to quarks and leptons. The second solution of the hierarchy
problem is based on supersymmetry. Supersymmetry is a symmetry that relates
bosons to fermions, and hence allows bosons to benefit from the chiral
protection afforded to fermion masses. It is also incredibly interesting because
it forces us to think about gravitational effects, and seems to be an intrinsic
component of string theory, our only successful theory of quantum gravity.
Generic supersymmetric extensions of the standard model, which allow for
parameters that break supersymmetry in the phenomenologically necessary way,
also have a variety of phenomenological problems. However, it is possible to
solve these problems in specific models. The key question seems to be to
understand the mechanism by which supersymmetry is broken. String theory
contains hints suggesting that supersymmetry breaking is much more constrained
than low-energy field theory would have us believe.

There is another resolution, or rather postponement, of the hierarchy problem
in a class of models that go under the rubric ``little Higgs''
\cite{banks-ref-119}. These models successfully postpone the hierarchy
discussion to the 100-TeV scale, and declare that, since that scale will be out
of experimental reach for the foreseeable future, we can safely ignore it. I
will leave the reader to decide on the value of these models for her/himself.

% BANKS-SOURCE: pdf=200 print=190 kind=subsection id=9.12.1
\subsubsection{Marginally relevant perturbations}

% BANKS-SOURCE: pdf=200 print=190 kind=prose
A remarkable theorem of Coleman, Gross, and Zee \cite{banks-ref-120,banks-ref-121,banks-ref-122}
shows that theories with marginally relevant parameters in four-dimensional
space-time always involve non-abelian gauge bosons. Non-abelian gauge couplings
are always marginally relevant unless the matter representation is too large.
Some of the other couplings in the theory may also turn out to be marginally
relevant, as a consequence of their interaction with the non-abelian bosons.
The result we have just proved suggests strongly that non-abelian gauge
interactions will be our only explanation of the large hierarchy

% BANKS-SOURCE: pdf=201 print=191 kind=source-page
in energy scales between the Planck mass, $10^{19}\,\mathrm{GeV}$, and the
typical scales of particle physics. This does indeed seem to be the explanation
of the scale of strong interactions. Notice that the phenomenon of confinement,
which we discussed in Chapter 8, involves a dimensionful parameter, the string
tension or energy per unit length of the QCD flux tube. If the QCD coupling
were not marginally relevant, we would expect this parameter to be at the
cut-off scale. The scale $\Lambda_{\mathrm{QCD}}$ is therefore sometimes called
the confinement scale of QCD, though experiment shows that the actual string
tension involves a scale roughly $2\pi$ larger than $\Lambda_{\mathrm{QCD}}$.

Actually, if we follow this philosophy stringently, we are led to either omit
fundamental scalar fields from our Lagrangian or insist on supersymmetry as a
fundamental principle. Supersymmetry relates scalar fields to Weyl fermions, in
such a way that chiral symmetries act on scalars. Thus, in supersymmetric
theories, scalar mass parameters, like fermion masses, behave like marginally
relevant operators, even though they have dimension $2$. They receive only
logarithmic renormalizations. Thus, it is consistent to have scalar mass
parameters of order $100\,\mathrm{GeV}$, even if the cut-off scale is much
higher than that. Supersymmetry is broken in the real world, and its virtues can
be preserved only if the superpartners of standard-model particles lie not too
far above the electro-weak scale.

This kind of solution of the hierarchy problem is technical in nature. It does
not yet explain the value of the electro-weak scale in the way that asymptotic
freedom explains the scale of strong interactions. In order to do that, we have
to make the breaking of the chiral symmetry that protects the Higgs mass into a
dynamical mechanism, which depends on some new marginally relevant parameter in
the Lagrangian of the world. This can be done either with or without
supersymmetry. In the non-supersymmetric solution, which is called technicolor
\cite{banks-ref-123}, the Higgs field is a fermion bilinear in a new QCD-like
sector. These models have severe phenomenological problems. There is a variety
of supersymmetric scenarios for a completely dynamical explanation of the
electro-weak scale. Many of them are roughly consistent with current
experimental data, but there are various causes for unease, and no uniquely
beautiful model has yet emerged.
